\chapter{Cupping}

\section*{Definition and Overview}
Cupping is a traditional therapy in which cups are placed on the skin and suction is applied, drawing the skin and underlying tissue upwards into the cup. It has been used for thousands of years in Chinese, Middle Eastern, and other traditional medical systems. Practitioners use cupping for pain, muscle tension, and a range of other conditions. The suction can be applied with heat or a pump, and cups may be left in place, moved over the skin, or, in wet cupping, combined with small skin incisions to draw a little blood. Evidence for the effectiveness of cupping is limited, and it should be used alongside, not instead of, conventional care.

\section*{Epidemiology}
Cupping is a popular complementary therapy worldwide, including in Australia, where it is offered by traditional medicine practitioners, acupuncturists, and some physiotherapists and massage therapists. People most commonly seek cupping for musculoskeletal pain, particularly back and neck pain, and for relaxation and sports recovery. Use has grown in recent years, partly through its visibility in elite sport and social media. Although widely used, cupping has limited high-quality research evidence, and its popularity reflects tradition and perceived benefit more than proven efficacy.

\section*{Aetiology and Pathophysiology}
Traditional theory explains cupping in terms of restoring the flow of energy (qi), while modern science offers alternative explanations.

\begin{itemize}
    \item \textbf{Traditional theory:} cupping is said to draw out toxins and restore balance in the body
    \item \textbf{Suction effects:} the suction lifts the skin and superficial muscles and increases local blood flow
    \item \textbf{Fascial effects:} some propose effects on connective tissue and pressure receptors
    \item \textbf{Pain modulation:} possible effects on pain pathways, similar to other physical therapies
    \item \textbf{Placebo and context:} expectation, relaxation, and the therapeutic relationship contribute to perceived benefit
    \item \textbf{The marks:} the characteristic circular bruises are caused by rupture of small blood vessels under the suction
\end{itemize}

\section*{Clinical Features}
Cupping is not a disease, so it has no symptoms of its own; it is used as a treatment.

\begin{itemize}
    \item \textbf{Common uses:} low back pain, neck and shoulder pain, muscle tension, and sports recovery
    \item \textbf{Other uses:} headache, fatigue, respiratory complaints, and general relaxation
    \item \textbf{The treatment:} cups are applied to the skin and left for several minutes
    \item \textbf{Sensations:} a pulling or tight sensation during treatment
    \item \textbf{After effects:} circular red or purple marks that fade over days to weeks
    \item \textbf{Side effects:} bruising, soreness, and, rarely, burns or skin infection
\end{itemize}

\section*{Differential Diagnosis}
Before using cupping, symptoms such as pain should be properly assessed, because serious conditions can be mistaken for minor problems.

\begin{itemize}
    \item Back pain due to fracture, infection, or cancer, or pressure on the spinal nerves
    \item Muscle and joint pain due to inflammatory arthritis needing treatment
    \item Heart or lung disease causing chest or arm pain
    \item Headaches caused by high blood pressure or other treatable conditions
    \item Skin conditions that make cupping unsafe, such as infection or fragile skin
    \item Bleeding disorders or blood-thinning medicines, which increase the risk of bruising and bleeding
\end{itemize}

\section*{When to Seek Medical Care}
See a doctor first to establish what is wrong before choosing cupping. Seek medical care if symptoms are severe, worsening, or accompanied by fever, weight loss, weakness, or numbness, or if they do not improve.

\textbf{Call 000 (or local emergency services) for:}
\begin{itemize}
    \item Chest pain, difficulty breathing, or sudden weakness
    \item Loss of consciousness, confusion, or a seizure
    \item A sudden severe headache, or new numbness or weakness of the face, arm, or leg
    \item Severe bleeding or signs of a serious skin infection after cupping
\end{itemize}

\section*{Diagnosis and Investigations}
Cupping itself requires no tests, but proper assessment of the condition being treated does.

\begin{itemize}
    \item \textbf{Medical assessment:} a doctor takes a history and examines the underlying condition
    \item \textbf{Investigations:} blood tests, X-rays, or scans may be arranged to confirm the cause of symptoms
    \item \textbf{Practitioner assessment:} a qualified practitioner takes their own history and checks for contraindications
    \item \textbf{Safety check:} discussion of medicines, pregnancy, skin conditions, and bleeding risk before treatment
\end{itemize}

\section*{Management}
\begin{itemize}
    \item \textbf{Qualified practitioner:} choose a registered or accredited practitioner with appropriate training
    \item \textbf{Treatment course:} usually a series of sessions, depending on the condition
    \item \textbf{Combined care:} best used alongside physiotherapy, exercise, and medicines, rather than instead of them
    \item \textbf{Exercise and self-management:} for conditions such as low back pain, staying active remains important
    \item \textbf{Review:} treatment should be reviewed and stopped if there is no benefit
    \item \textbf{Conventional follow-up:} continue to see your doctor and attend planned tests or reviews
\end{itemize}

\section*{Complications}
Cupping is generally safe in trained hands, but complications can occur.

\begin{itemize}
    \item Bruising and circular marks, which are expected
    \item Soreness and discomfort after treatment
    \item Burns, particularly with fire cupping
    \item Skin infection, especially with wet cupping or unsterile equipment
    \item Bleeding or worsening of anaemia with frequent wet cupping
    \item Dizziness or fainting during treatment
    \item The risk of relying on cupping while a serious condition goes undiagnosed
\end{itemize}

\section*{Prognosis and Outcomes}
Evidence for cupping is limited, and any benefit for pain appears modest and is likely due in part to expectation and relaxation. Many people report temporary relief of muscle tension and a sense of wellbeing, but cupping is not a cure for underlying disease. It is safest and most useful as a complementary therapy alongside conventional care. Anyone with persistent or serious symptoms should have a proper medical diagnosis and treatment rather than relying on cupping alone.

\section*{Prevention and Self-Care}
\begin{itemize}
    \item Seek a proper medical diagnosis before using cupping for a symptom
    \item Choose a qualified practitioner who uses clean, hygienic equipment
    \item Tell the practitioner about medicines, especially blood thinners, skin conditions, and pregnancy
    \item Avoid cupping over broken skin, varicose veins, or inflamed areas
    \item Do not delay conventional treatment for symptoms that could be serious
    \item Review the treatment after a few sessions and stop if it is not helping
\end{itemize}

\section*{Sources and Further Reading}
The following reputable sources provide further information for healthcare students and patients seeking reliable, current advice about cupping and complementary medicine.

\begin{itemize}
    \item NHS. \textit{Cupping and complementary therapies.} https://www.nhs.uk/conditions/
    \item National Center for Complementary and Integrative Health. \textit{Traditional Chinese medicine and cupping.} https://www.nccih.nih.gov/
    \item Australian Health Practitioner Regulation Agency (AHPRA). \textit{Chinese medicine practitioner registration.} https://www.ahpra.gov.au/
    \item World Health Organization. \textit{Traditional and complementary medicine.} https://www.who.int/
    \item Healthdirect Australia. \textit{Complementary therapies.} https://www.healthdirect.gov.au/
    \item Mayo Clinic. \textit{Complementary and alternative medicine.} https://www.mayoclinic.org/
\end{itemize}
